\documentclass[11pt,a4paper]{article}

\usepackage[T1]{fontenc}
\usepackage[utf8]{inputenc}
\usepackage{lmodern}
\usepackage{microtype}
\usepackage{amsmath,amssymb}
\usepackage{booktabs}
\usepackage{tabularx}
\usepackage{longtable}
\usepackage{array}
\usepackage{xcolor}
\usepackage{geometry}
\usepackage{hyperref}
\usepackage{setspace}

\geometry{margin=27mm}
\setstretch{1.04}
\hypersetup{
  colorlinks=true,
  linkcolor=blue!50!black,
  citecolor=blue!50!black,
  urlcolor=blue!50!black,
  pdftitle={Comparison of Local-Module and Enlarged-Subnetwork Identification Methods},
  pdfauthor={Anders Hansson}
}

\newcommand{\grank}{\operatorname{grank}}
\newcommand{\specmat}[1]{\Phi_{#1}}
\newcommand{\rank}{\operatorname{rank}}
\newcommand{\fiber}{\mathcal F}
\newcommand{\Aset}{\mathcal A}
\newcommand{\Dset}{\mathcal D}
\newcommand{\Rset}{\mathcal R}

\title{Comparison of Local-Module and Enlarged-Subnetwork Identification Methods\\
\large Plant Modules, Measurements, Informativity, and Identifiability}
\author{Anders Hansson}
\date{}

\begin{document}
\maketitle

\begin{abstract}
This note compares the prediction methods of Van den Hof et al. (2013), Dankers et al. (2016), and Weerts et al. (2020) with the enlarged-subnetwork method of da Mata, Hansson, and Andersen.  It distinguishes three questions: whether prediction errors have zero cross-moments with the permitted data, whether the experiment distinguishes candidate predictors, and whether the original plant modules can be recovered from measured transfer functions.  The contribution returning from the surrounding network must be reconstructible from the permitted measurement history.  Two constructed examples demonstrate algebraic recovery with only the target input and output measured, or with neither measured.  These calculations alone do not establish consistent estimation.  Conclusions are conditional where the topology, experiment, or predictor is unspecified.  No minimum enlargement or weaker topology-only algorithm is established.
\end{abstract}

\section{Purpose and scope}

The purpose is to examine the enlarged-subnetwork method on the displayed reference configurations, rather than only on its own numerical example.  Two final constructions illustrate algebraic recovery with reduced internal measurements.  Comparisons concern the particular predictor configurations displayed below, not a proof of the minimum measurement requirements of every method in the references.

The note concerns recovery of plant transfer functions, not disturbance filters.  Definitions precede the assessment criteria, reference configurations, and constructed examples.  The final tables collect the comparisons.

\section{Definitions and assessment criteria}
\label{sec:framework}

\subsection{Network notation and predictor information}
\label{sec:notation}

A plant module is a transfer function connecting an input signal to an output signal.  A plant model consists of specified signal equations and candidate plant modules; a noise model specifies filters and covariance matrices for the additive disturbances.  A stochastic model combines these plant and noise models.  A target module is selected for recovery; other unknown plant modules are nuisance modules.  The target endpoints are its input and output signals.  Write $w_i$ for node signal $i$, $r_i$ for an external excitation, $v_i$ for an additive disturbance, and $G_{ji}$ for the module from $w_i$ to $w_j$.  Subscript $k$ denotes discrete time.  The operator $q$ is the forward time shift, so $q^{-1}$ is a one-sample delay.  The directed network topology records the signal-to-signal connections, not their numerical transfer functions.

Let $\Aset$ denote the selected set of plant modules to be estimated and $\mathcal B$ the complementary part.  Enlargement moves original modules from $\mathcal B$ into $\Aset$.  Write $y^{\Aset}$ and $u^{\Aset}$ for the selected module-output and module-input vectors.  Let $x$ be the measured predictor-output vector; in the supplied manuscript it is a measured part of $(y^{\Aset},u^{\Aset})$.  Let $r^{\mathrm{exo}}$ collect the available measured exogenous signals and $\bar y^B$ the measured boundary signals entering $\Aset$ from $\mathcal B$.  The augmented input is $\bar r=((r^{\mathrm{exo}})^T,(\bar y^B)^T)^T$.  The predictor information consists of $x_j$ for $j<k$ and the available augmented-input values through time $k$.  A one-step predictor $\hat x_k$ estimates $x_k$ from this information; its prediction error is $\varepsilon_k=x_k-\hat x_k$.  A predictor model is the set of candidate one-step predictors obtained by varying the plant and noise parameters.

Let $\theta$ parameterize the unknown modules, $\theta^0$ be the true parameter, and $\theta_{\Aset}$ denote the parameters of the selected original modules.  Consistency means convergence of the estimates to the true modules as the amount of data increases.  Each network has its own signals and modules; repeated node indices do not identify signals across networks.  Each configuration specializes $\Aset$, $x$, $s$, and $u^{\Aset}$ according to these definitions.

For plant calculations, let $s$ collect the selected measured exogenous and incoming boundary signals.  Define the plant transfer matrices by suppressing disturbances and treating $s$ as the input of the subnetwork equations:
\begin{equation}
  u^{\Aset}=T_{us}(q,\theta)s,
  \qquad x=T_{xs}(q,\theta)s.
  \label{eq:plant-transfers}
\end{equation}
The measured transfer matrix is $T_{xs}(q,\theta)$; the measured transfer map is the parameter-to-matrix function $\theta\mapsto T_{xs}(q,\theta)$.  Subscripts on $T$ name its output and input signals; for example, $T_{w_jr_i}$ is a scalar transfer.  Endogenous signals depend on network disturbances or outputs; exogenous signals are supplied independently of them in the experiments considered here.  When $s$ contains endogenous boundary signals, the plant transfer matrices in \eqref{eq:plant-transfers} do not describe a regression on independent excitations.  They are not the noise-filtered one-step predictor.

\subsection{Genericity, rank, and spectral notation}
\label{sec:generic-definitions}

Let $\Theta_{\mathrm{ad}}\subseteq\mathbb R^p$ describe a specified finite-dimensional admissible parameter class.  A statement is universal if it holds for every $\theta\in\Theta_{\mathrm{ad}}$.  It holds almost everywhere if it holds outside an exceptional set $N$ of Lebesgue measure zero.  For a property $\mathcal P(\theta)$ of the parameterized network, the two guarantees are respectively
\[
  \mathcal P(\theta)\quad\text{for every }\theta\in\Theta_{\mathrm{ad}}
\]
and
\begin{equation}
  \mathcal P(\theta)\quad\text{for every }
  \theta\in\Theta_{\mathrm{ad}}\setminus N,
  \qquad N\text{ has Lebesgue measure zero}.
  \label{eq:generic-boundary-condition}
\end{equation}
An algebraic exceptional set is defined by polynomial equations in the parameters; it is proper when those equations do not hold throughout the admissible parameter class.  Here, algebraic genericity means validity outside such a set of measure zero.  A certificate is an argument establishing a property universally or generically; a sufficient certificate need not cover every case where the property holds.  A certificate is strictly conservative if it rejects an admissible case where the property holds.  Minimum enlargement minimizes a stated cost over stated choices of added modules.

For a rational transfer matrix, $\rank$ denotes its normal rank, that is, its rank over the rational-function field.  The generic rank $\grank$ is the normal rank attained away from exceptional parameter values in the specified class.  For a stationary vector signal $a$, $\specmat{a}(\omega)$ denotes its spectral-density matrix at frequency $\omega$.  The notation $\specmat{a}(\omega)\succ0$ means that this matrix is positive definite; scalar positivity is written $\specmat{a}(\omega)>0$.  Spectral conditions are understood almost everywhere in frequency.  The dimension of a vector is written $\dim$, and $I_n$ denotes the $n\times n$ identity matrix.

\subsection{Orthogonality and boundary reconstructibility}
\label{sec:orthogonality-definitions}

A rational transfer function in $q$ is proper if its numerator degree does not exceed its denominator degree; it is strictly proper if that degree is smaller.  Strict properness excludes instantaneous input-to-output transmission.  A stable transfer has its poles strictly inside the unit circle; closed-loop stability concerns the interconnected network.  A causal reconstruction uses no future signal values.  Reconstruction from measured history means exact recovery at time $k$ from past measured outputs and exogenous signals available through time $k$.  An innovation is a new unpredictable noise component relative to the specified past information.  Disturbance separation means diagonal disturbance filters and diagonal innovation covariance, separating the innovations of the selected and complementary parts.

Prediction-error orthogonality means the following zero cross-moment relations:
\begin{equation}
  \mathbb E[\varepsilon_kx_j^T]=0\quad(j<k),
  \qquad
  \mathbb E[\varepsilon_k\bar r_j^T]=0\quad(j\leq k).
  \label{eq:orthogonality}
\end{equation}
Let $b^{\Aset}$ denote the component of $\bar y^B$ generated by signals from $\Aset$, as defined in the supplied manuscript.  The boundary-information condition requires $b^{\Aset}_k$ to be causally reconstructible from $\{x_j:j<k\}$ and the measured exogenous signals through time $k$.  The manuscript uses three conditions:
\begin{enumerate}
  \item[(O1)] the plant transfer matrix is strictly proper;
  \item[(O2)] disturbance separation holds; and
  \item[(O3)] the boundary-information condition holds.
\end{enumerate}
Conditions (O1)--(O2) concern the plant and noise models; (O3) concerns the selected modules and predictor information.  Under these conditions, the orthogonality argument in da Mata, Hansson, and Andersen~\cite{DaMataHanssonAndersenSubnetwork} establishes \eqref{eq:orthogonality}.  Strict properness prevents instantaneous transmission of current innovations to returning boundary signals; disturbance separation removes correlation with innovations in the complementary part.
Two polynomials are coprime if they share no nonconstant factor.  The predictor regularity conditions are coprimeness of the plant polynomials and absence of noise-polynomial zeros on the unit circle, as required in the supplied consistency theorem.  They are separate from (O1)--(O3).

Reconstructibility from the history of another measured signal outside this specified model is not sufficient to invoke the theorem unless that signal's inclusion in the predictor and noise model has also been justified.

\subsection{Excitation and predictor informativity}
\label{sec:excitation-definitions}

Excitation refers to variation in the signals used as model inputs.  Informativity concerns whether the experimental data can distinguish the candidate predictors in the specified model class.  Define the actual predictor-data vector as $z=(x^T,\bar r^T)^T$.  A plant-input excitation rank test used below is
\begin{equation}
  \specmat{s}(\omega) \succ 0
  \quad\text{and}\quad
  \grank T_{u s}(q,\theta)=\dim u^{\Aset}.
  \label{eq:informativity}
\end{equation}
For exogenous $s$ independent of the disturbances, the externally generated plant-input spectrum is
\begin{equation}
  T_{us}(e^{\mathrm i\omega},\theta)\specmat{s}(\omega)
  T_{us}(e^{\mathrm i\omega},\theta)^* ,
  \label{eq:external-input-spectrum}
\end{equation}
where $^*$ denotes conjugate transpose.  Under \eqref{eq:informativity}, at nonexceptional parameters and almost every frequency, $T_{us}^*a\ne0$ for every nonzero vector $a$; the corresponding quadratic form in \eqref{eq:external-input-spectrum} is positive.  Endogenous boundary signals require a separate argument: an identity or selection matrix does not establish independence from disturbances.  The plant-input excitation rank test is not a general certificate of predictor informativity.  Hansson, da Mata, and Andersen~\cite{HanssonDaMataAndersen2026} give a sufficient spectral condition involving $\specmat{z}(\omega)\succ0$, under their predictor assumptions.  Neither spectral condition establishes \eqref{eq:orthogonality}.

\subsection{Plant-module identifiability and fibers}
\label{sec:fiber-definitions}

Plant-module identifiability concerns recovery of the original modules from the exact measured transfer matrix.  Define the fiber at the true parameter $\theta^0$ by
\begin{equation}
  \fiber(\theta^0)
  :=\{\theta\in\Theta_{\mathrm{ad}}:T_{xs}(q,\theta)=T_{xs}(q,\theta^0)\}.
  \label{eq:fiber}
\end{equation}
Local identifiability means uniqueness near the true parameter; global identifiability means uniqueness over the admissible parameter class.  A parameterization is nonredundant when distinct parameters describe distinct module tuples.  For a specified nonredundant algebraic parameterization, the dimension of a fiber counts its locally independent parameter directions; dimension zero means isolated solutions.  A zero-dimensional generic fiber establishes generic local identifiability; a generically singleton fiber establishes generic global identifiability~\cite{HanssonDaMataAndersen2026}.  A target module is globally identifiable when it is constant on the fiber in \eqref{eq:fiber}, even if nuisance modules vary.  A positive-dimensional fiber therefore does not by itself prove target nonidentifiability.  Orthogonality, predictor informativity, and plant-module identifiability require separate checks.

Products and feedback denominators such as $G_aG_b$ or $(1-G_aG_b)^{-1}$ do not themselves imply nonidentifiability.  The question is whether the measured transfer matrix $T_{xs}$ provides enough independent relations to make the relevant fiber zero-dimensional.

\subsection{Topology-based and generic boundary-information certificates}
\label{sec:generic-boundary-certificate}

Departures are signals leaving $\Aset$ and entering $\mathcal B$; let $\mathcal D_{\mathrm{out}}$ collect them.  A returning path is a directed path through $\mathcal B$ from a departure to a measured incoming boundary signal.  Define
\[
  \mathcal D_{\mathrm{return}}
  :=\{d\in\mathcal D_{\mathrm{out}}:\text{a returning path starts at }d\}.
\]
The departure certificate requires every signal in $\mathcal D_{\mathrm{return}}$ to be measured or causally reconstructible within the predictor information, with a strictly proper transfer on every returning path.  This is the sufficient certificate for (O3) supported by the supplied manuscript; it does not require measurement of departures outside $\mathcal D_{\mathrm{return}}$.  Enlargement moves modules into $\Aset$ and the departure certificate must then be checked at the new boundary.  If the complete relevant network is included, no returning boundary remains and (O3) is vacuous.  Failure of the departure certificate does not prove failure of (O3), whose requirement concerns only $b^{\Aset}$.  Neither minimum enlargement nor strict conservatism of the departure certificate is established here.

\paragraph{Applying universal and generic guarantees.}
Specialize the property $\mathcal P(\theta)$ to (O3), with $\Theta_{\mathrm{ad}}$ satisfying the standing stability, strict-properness, and disturbance assumptions.  A certificate may establish the generic guarantee in \eqref{eq:generic-boundary-condition} rather than universal validity.  A related distinction appears in results on recovering functions of the system state from measurements~\cite{ZhangFernandoDarouachSFO}; those results are not a theorem for (O3).

For $\theta^0\notin N$, the certificate establishes (O3) at the true parameter.  Under (O1)--(O2) and the remaining predictor assumptions, the supplied orthogonality argument then establishes \eqref{eq:orthogonality}.  This gives generic orthogonality, not a guarantee for $\theta^0\in N$.  Reconstruction may depend on the unknown true modules: neither universal nor generic validity requires one filter to work for every parameter value or numerical knowledge of its coefficients when checking the certificate.

Topology and strict properness do not specify $\Theta_{\mathrm{ad}}$ or a measure over transfer functions of unrestricted order.  A generic claim must state the parameter dependencies and order or delay restrictions used in its proof.  This does not require numerical knowledge of the modules.  Disturbance separation and the specified predictor information remain essential.

\paragraph{Algebraic screening is not a causal certificate.}
For a specified rational model class, let $\xi$ collect the primitive excitation and disturbance sources of the network.  Define the reconstruction-data vector $z^{\mathrm{rec}}$ and its source-to-data transfer $T_{\mathrm{rec}}$, and let $T_b$ be the source-to-returning-component transfer:
\[
  z^{\mathrm{rec}}_k=\begin{bmatrix}x_{k-1}\\r^{\mathrm{exo}}_k\end{bmatrix},
  \qquad
  z^{\mathrm{rec}}=T_{\mathrm{rec}}(q,\theta)\xi,
  \qquad
  b^{\Aset}=T_b(q,\theta)\xi,
\]
The delay distinguishes $z^{\mathrm{rec}}$ from the predictor-data vector $z$: only strictly past $x$ is available.  A reconstruction filter $L_{\mathrm{rec}}(q,\theta)$ satisfying
\begin{equation}
  T_b(q,\theta)=L_{\mathrm{rec}}(q,\theta)T_{\mathrm{rec}}(q,\theta)
  \label{eq:boundary-factorization}
\end{equation}
certifies (O3) only if it defines a causal, well-defined reconstruction on the signal class in the theorem.  All primitive-source columns must be retained: a reconstruction must not inadvertently retain an unwanted disturbance contribution from a measured signal.

Ignoring causality and the permitted signal histories, each row of $T_b$ is a rational linear combination of the rows of $T_{\mathrm{rec}}$ generically exactly when
\begin{equation}
  \grank\!\begin{bmatrix}T_{\mathrm{rec}}\\T_b\end{bmatrix}
  =\grank T_{\mathrm{rec}}.
  \label{eq:boundary-rank-screen}
\end{equation}
Condition \eqref{eq:boundary-rank-screen} tests existence of a rational factor in \eqref{eq:boundary-factorization}; shared module parameters must be preserved.  A factor exists exactly when adding the rows of $T_b$ does not enlarge the linear span of the rows of $T_{\mathrm{rec}}$ over rational functions.  This does not prove causality or stability.  For example, the strictly proper transfer $q^{-1}$ has the inverse $q$, which uses a future sample.  The rank test therefore needs a separate causal reconstruction argument.  It does not establish a weaker topology-only algorithm for (O3).

\subsection{Local-model terminology and comparison order}
\label{sec:comparison-protocol}

A single-input single-output (SISO) configuration has one plant input and one output; a multiple-input single-output (MISO) configuration has several plant inputs and one output.  A direct local configuration uses measured node inputs in a selected node equation.  A two-stage configuration reconstructs input components from measured external excitation before using them in the local plant model.  Immersion eliminates omitted nodes from the network equations.  An indirect observation is a measured signal related to an unmeasured node through plant modules.  Abstraction transforms the network equations using retained signals or indirect observations.  Target invariance means that the transformed target coefficient equals the original target module; a transformation satisfying this requirement is target-preserving.  Path- and loop-blocking conditions require retained signals to intercept alternative paths or loops relevant to target invariance.  These definitions do not guarantee consistency.

Confounding refers here to an unmeasured disturbance source affecting both a model input and its output.  This concept is distinct from the boundary-information condition used in the enlarged-subnetwork argument.

For each example, the analysis first states the network, target, and selected measurements, then records:
\begin{enumerate}
  \item the open-loop or abstracted transfer functions occurring in the predictor;
  \item the measured output, node, boundary, and external signals;
  \item satisfaction of the strict-properness, disturbance-separation, and boundary-information conditions used to obtain prediction-error orthogonality;
  \item excitation conditions, distinguished from predictor informativity; and
  \item generic local or global identifiability of the target modules from the measured transfer matrix.
\end{enumerate}
The comparison concerns plant-module identifiability.  Noise-model identifiability is not claimed.  The plant-input excitation tests do not by themselves establish predictor informativity, and satisfaction of an excitation rank condition is never used as a substitute for the separate boundary-information condition.

\section{Van den Hof et al. (2013)}

\subsection{Common five-node network}

The numbered examples use the network
\begin{align}
  w_1 &= G_{15}w_5+r_1, \label{eq:vdh1}\\
  w_2 &= G_{21}w_1+G_{23}w_3+v_2, \label{eq:vdh2}\\
  w_3 &= G_{32}w_2+v_3, \label{eq:vdh3}\\
  w_4 &= w_3+G_{45}w_5, \label{eq:vdh4}\\
  w_5 &= G_{54}w_4+v_5, \label{eq:vdh5}
\end{align}
where $G_{43}=1$.  The paper treats $G_{15}$, $G_{32}$, $G_{54}$, and $G_{45}$ with SISO predictors and treats $G_{21}$ and $G_{23}$ jointly with a MISO predictor.  The known static link $G_{43}=1$ is not strictly proper.  To apply a theorem requiring strictly proper plant modules, this link must enter the known interconnection equations, with the resulting plant and noise models checked separately.  The comparisons check (O3) and the plant algebra, not this complete predictor construction.

\subsection{Example 1: direct SISO predictors}

For the four SISO targets, both methods use the same plant relation.  Each configuration selects the single target equation as the $\Aset$-part.

\begin{table}[htbp]
\centering
\caption{Van den Hof et al. Example 1: direct SISO configurations.}
\label{tab:vdh-siso}
\begin{tabularx}{\textwidth}{@{}lX X X@{}}
\toprule
Target & Plant relation used by both methods & Measured signals & Modules in predictor \\
\midrule
$G_{15}$ & $w_1=G_{15}w_5+r_1$ & $w_1,w_5,r_1$ & $G_{15}$ \\
$G_{32}$ & $w_3=G_{32}w_2+v_3$ & $w_3,w_2$ & $G_{32}$ \\
$G_{54}$ & $w_5=G_{54}w_4+v_5$ & $w_5,w_4$ & $G_{54}$ \\
$G_{45}$ & $w_4-w_3=G_{45}w_5$ & $w_4,w_3,w_5$ & $G_{45}$ \\
\bottomrule
\end{tabularx}
\end{table}

For the enlarged-subnetwork interpretation, the orthogonality condition is satisfied by the departure certificate.  For $G_{15}$, $G_{32}$, and $G_{54}$, the output of the scalar $\Aset$-part is respectively $w_1$, $w_3$, and $w_5$, and that signal is measured before it can propagate through the surrounding network and return as a predictor input.  For $G_{45}$, the module output is causally reconstructible as $w_4-w_3$ from the measured signals.  Thus the component of every returning boundary signal generated by the scalar $\Aset$-part is reconstructible from measured history.  Under the strict-properness and disturbance-separation assumptions, \eqref{eq:orthogonality} follows.

Each displayed plant equation has one scalar input.  If that input is a measured boundary signal, the plant-input transfer matrix is $T_{us}=1$, and positivity of its spectrum supplies scalar plant-input excitation.  For a component reconstructed from $r_1$, a nonzero transfer from $r_1$ and positive excitation spectrum supply a nonzero reconstructed-input spectrum.  Neither statement alone verifies all assumptions for the stochastic predictor.

The measured transfer matrix exposes the target module directly.  For example,
\[
  T_{w_1w_5}=G_{15},
  \qquad
  T_{w_3w_2}=G_{32}.
\]
The corresponding generic fibers are singletons.  Thus the enlarged-subnetwork method gives no measurement or identifiability advantage in these SISO cases.

\subsection{Example 1: direct MISO predictor for $G_{21}$ and $G_{23}$}

Both methods use
\begin{equation}
  w_2=G_{21}w_1+G_{23}w_3+v_2,
  \label{eq:vdh-miso}
\end{equation}
with measured signals $w_1,w_2,w_3$.  The predictor-input vector is
\[
  u^{\Aset}=\begin{bmatrix}w_1&w_3\end{bmatrix}^{\!T}.
\]
The departure signal $w_2$ is measured.  The departure certificate therefore establishes (O3) for this configuration.  Under (O1)--(O2), the orthogonality argument applies to the measured incoming boundary signals.
If the two inputs are to be reconstructed from $r_1$ alone, then
\[
  T_{u r_1}\in\mathbb R(q)^{2\times 1},
  \qquad
  \rank T_{u r_1}\leq 1,
\]
so the plant-input excitation rank test does not pass for two reconstructed inputs.  This does not prove failure of predictor informativity.

For the enlarged-subnetwork formulation, $w_1$ and $w_3$ are measured boundary signals.  With
\[
  s=\begin{bmatrix}w_1&w_3\end{bmatrix}^{\!T},
  \qquad T_{us}=I_2,
\]
condition \eqref{eq:informativity} reduces to
\begin{equation}
  \specmat{[w_1,w_3]^T}(\omega)\succ0.
  \label{eq:vdh-boundary-info}
\end{equation}
The measured transfer matrix is
\begin{equation}
  T_{w_2,[w_1,w_3]}
   =\begin{bmatrix}G_{21}&G_{23}\end{bmatrix}.
  \label{eq:vdh-target-matrix}
\end{equation}
Its entries in \eqref{eq:vdh-target-matrix} determine both plant modules.  The configurations use the same node sensors.  Condition \eqref{eq:vdh-boundary-info} checks plant-input excitation, not predictor informativity.

\subsection{Example 2: SISO two-stage identification}

The referenced method replaces each scalar measured predictor input $w_i$ by its component $w_i^{(r_1)}$ reconstructed from $r_1$.  It is applied to $G_{15}$, $G_{32}$, $G_{54}$, and $G_{45}$.  For every scalar target, the rank part of \eqref{eq:informativity} is satisfied when the transfer $T_{w_ir_1}$ is nonzero.  The excitation spectrum must additionally satisfy $\specmat{r_1}(\omega)>0$.

The enlarged-subnetwork configuration instead uses $w_i$ directly.  The condition $\specmat{w_i}(\omega)>0$, with $T_{us}=1$, checks scalar plant-input excitation.  In both displayed plant-level formulations, the target module is exposed directly.  The configurations use the same node sensors; the enlarged-subnetwork configuration does not project the input onto $r_1$.

The orthogonality conclusion is unchanged from the direct SISO configurations: the output of each scalar $\Aset$-part is measured, and the $\Aset$-generated component of a returning boundary signal is reconstructible from that measured output history.  Projection onto $r_1$ concerns informativity and confounding in the referenced two-stage construction; it is not the boundary-information condition used by the enlarged-subnetwork method.

\subsection{Example 3: MISO two-stage identification}

The referenced predictor is
\begin{equation}
  \widehat w_2
  =G_{21}w_1^{(r_1)}+G_{23}w_3^{(r_1)}.
  \label{eq:vdh-twostage-miso}
\end{equation}
With only scalar $r_1$, the transfer matrix from the external excitation to the two reconstructed predictor inputs is $2\times1$ and cannot have full row rank two.  Therefore \eqref{eq:informativity} does not certify this experiment.  A second independent excitation can make the dimensions compatible with the sufficient rank condition.

The enlarged-subnetwork configuration again has $T_{us}=I_2$ for $s=(w_1,w_3)^T$.  Condition \eqref{eq:vdh-boundary-info} checks plant-input excitation, and the measured transfer matrix exposes $G_{21}$ and $G_{23}$ separately.  It does not, on its own, certify the complete predictor.  This configuration uses the same node sensors without reconstructing both inputs from the scalar excitation.

Because $w_2$ is measured, the same structural argument as in the direct MISO case establishes the boundary-information condition.  The failure of the $2\times1$ rank test for the referenced two-stage construction is therefore an informativity issue and not a failure of prediction-error orthogonality for the enlarged-subnetwork formulation.

\subsection{Example 4: reconstructed disturbance}

After removal of the outer-loop module $G_{15}$, the referenced method can reconstruct
\begin{equation}
  v_3=w_3-G_{32}w_2
  \label{eq:reconstructed-v3}
\end{equation}
when $G_{32}$ is known.  The reconstructed disturbance acts as an excitation for identification of $G_{23}$.  For a two-input identification of $G_{21}$ and $G_{23}$, the measurable excitation vector can be taken as $(r_1,v_3)^T$.  A sufficient condition is
\begin{equation}
  \specmat{[r_1,v_3]^T}(\omega)\succ0,
  \qquad
  \grank T_{[w_1,w_3]^T,[r_1,v_3]^T}=2.
  \label{eq:vdh-v3-info}
\end{equation}

The enlarged-subnetwork configuration takes $x=w_2$, $s=(w_1,w_3)^T$, and $u^{\Aset}=s$.  Thus $T_{us}=I_2$, and \eqref{eq:vdh-boundary-info} checks plant-input excitation.  The measured transfer matrix is unchanged from \eqref{eq:vdh-target-matrix}.
This plant-level formulation uses neither reconstruction of $v_3$ nor prior knowledge of $G_{32}$; the node measurements are unchanged.  A consistency comparison still requires the appropriate noise model and predictor assumptions.

Here again $w_2$ is the measured output of the $\Aset$-part, so its contribution to boundary signals returning through the surrounding network is reconstructible from measured history.  Reconstruction of $v_3$ in the referenced method is therefore not needed to establish the orthogonality condition for the enlarged-subnetwork predictor.

\section{Dankers et al. (2016)}

Select $G_{ji}$ as the target module.  Let $\Dset_j$ be the retained predictor-input node indices and $w_{\Dset_j}$ their signal vector.  Write $\breve G_{jk}$ for the plant coefficients after immersion and $\breve v_j$ for the resulting disturbance.  Dankers et al. require target invariance, $\breve G_{ji}=G_{ji}$, under their path- and loop-blocking conditions.  The resulting plant relation is
\begin{equation}
  w_j
  =G_{ji}w_i
  +\sum_{k\in\Dset_j\setminus\{i\}}\breve G_{jk}w_k
  +\breve v_j,
  \label{eq:dankers-predictor}
\end{equation}

\subsection{Informativity of the referenced predictor}

For a direct experiment, a sufficient condition is
\begin{equation}
  \specmat{w_{\Dset_j}}(\omega)\succ0.
  \label{eq:dankers-direct-info}
\end{equation}
Let $\Rset$ be the selected set of external excitation indices and let $r_{\Rset}$ stack those excitation signals.  For a two-stage construction based on this vector, the corresponding sufficient condition is
\begin{equation}
  \specmat{r_{\Rset}}(\omega)\succ0,
  \qquad
  \grank T_{w_{\Dset_j},r_{\Rset}}=|\Dset_j|.
  \label{eq:dankers-twostage-info}
\end{equation}
Thus the rank test requires enough independent excitation directions to span the selected predictor-input vector.

\subsection{Identifiability of the referenced predictor}

Because the target is invariant and appears as a separate coefficient of $w_i$ in \eqref{eq:dankers-predictor}, the target is directly identifiable once the predictor is informative, subject to a nonredundant parameterization.  The nuisance coefficients $\breve G_{jk}$ need not equal individual original modules; they can be composite transfer functions generated by immersion.

\subsection{Applying the enlarged-subnetwork method}

An enlarged-subnetwork construction can include original modules removed by immersion and use retained node signals as measured boundary inputs.  Identity or selection columns in $T_{us}$ then concern plant-input excitation only.  They do not certify informativity or consistency of the complete predictor.

For every selected $\Aset$-part, (O3) must be verified using the actual predictor information set.  Measuring or reconstructing the relevant departure signals is a sufficient route.  Enlarging the model is another sufficient remedy; necessity or minimality of that remedy is not established here.

The enlarged construction requires analysis of the measured transfer map for the selected original modules, rather than recovery of $G_{ji}$ as an isolated immersed coefficient.  Apply the fiber criterion in Section~\ref{sec:fiber-definitions}; target identifiability requires checking whether the target varies within the fiber.

The exact edge equations of the numbered running example in Dankers et al. are not specified in this note.  Neither the boundary-information condition nor the algebraic fiber is therefore verified for that example.  The enlarged-subnetwork conclusions remain conditional on verifying (O3) for the specified predictor and on checking target identifiability for its measured transfer matrix.

\section{Weerts et al. (2020)}

\subsection{Four-node abstraction example}

The network is
\begin{align}
  w_1&=G_{12}w_2+G_{13}w_3+G_{14}w_4+r_1+v_1,\\
  w_2&=G_{24}w_4+r_2+v_2,\\
  w_3&=r_3+v_3,\\
  w_4&=G_{41}w_1+r_4+v_4,
  \label{eq:weerts-four-node}
\end{align}
with target $G_{13}$.

\subsubsection{Referenced abstraction}

The measured node signals are $w_1,w_2,w_3$.  Define $\widetilde v_1^{\mathrm{ref}}=v_1-G_{14}G_{24}^{-1}v_2$.  The second relation in \eqref{eq:weerts-four-node} gives $w_4=G_{24}^{-1}(w_2-r_2-v_2)$.  Substitution into its first relation gives
\begin{equation}
  w_1=
  \left(G_{12}+G_{14}G_{24}^{-1}\right)w_2
  +G_{13}w_3
  +r_1-G_{14}G_{24}^{-1}r_2
  +\widetilde v_1^{\mathrm{ref}}.
  \label{eq:weerts-indirect}
\end{equation}
The target remains the coefficient of $w_3$, whereas the nuisance coefficient is composite.  A sufficient direct-informativity condition is
\begin{equation}
  \specmat{[w_2,w_3]^T}(\omega)\succ0.
  \label{eq:weerts-reference-info}
\end{equation}
Algebraically, the target $G_{13}$ remains a separate coefficient; recovery of $G_{12}$, $G_{14}$, and $G_{24}$ individually is not required.  This elimination is not itself a causal-predictor construction: for strictly proper $G_{24}$, its inverse need not be causal or stable.  Predictor implementation and consistency require the regularity and noise-treatment assumptions of the referenced method.

\subsubsection{Enlarged-subnetwork construction}

Select the modules in the first, second, and fourth relations of \eqref{eq:weerts-four-node}.  Take
\[
  x=\begin{bmatrix}w_1&w_2\end{bmatrix}^{\!T},
  \qquad
  s=\begin{bmatrix}w_3&r_1&r_2&r_4\end{bmatrix}^{\!T}.
\]
The selected original modules are
\[
  \Aset=\{G_{12},G_{13},G_{14},G_{24},G_{41}\}.
\]
The only remaining boundary node in the displayed network is $w_3=r_3+v_3$.  It receives no signal from the $\Aset$-part, so its $\Aset$-generated component is identically zero.  The boundary-information condition therefore holds without a further enlargement, and the measured use of $w_3$ is compatible with \eqref{eq:orthogonality}.

Define $D_{\mathrm W}=1-G_{14}G_{41}$, $\widetilde v_1^{\mathrm{enl}}=(v_1+G_{14}v_4)/D_{\mathrm W}$, and $\widetilde v_2^{\mathrm{enl}}=v_2+G_{24}v_4$.  Substitute $w_4=G_{41}w_1+r_4+v_4$ from \eqref{eq:weerts-four-node} into its first two relations and collect the $w_1$ terms:
\begin{align}
  w_1
  &=\frac{G_{12}}{D_{\mathrm W}}w_2
    +\frac{G_{13}}{D_{\mathrm W}}w_3
    +\frac{1}{D_{\mathrm W}}r_1
    +\frac{G_{14}}{D_{\mathrm W}}r_4
    +\widetilde v_1^{\mathrm{enl}},\\
  w_2
  &=G_{24}G_{41}w_1+G_{24}r_4+r_2+\widetilde v_2^{\mathrm{enl}}.
  \label{eq:weerts-enlarged-relations}
\end{align}
The coefficients in \eqref{eq:weerts-enlarged-relations} are not the closed-loop transfer entries.  Define
\[
  C_{\mathrm W}=1-G_{41}(G_{14}+G_{12}G_{24}).
\]
For the plant transfer calculation, suppress disturbances and substitute the second relation of \eqref{eq:weerts-enlarged-relations} into the first:
\begin{equation}
  C_{\mathrm W}w_1
  =G_{13}w_3+r_1+G_{12}r_2+(G_{14}+G_{12}G_{24})r_4.
  \label{eq:weerts-plant-elimination}
\end{equation}
Divide \eqref{eq:weerts-plant-elimination} by $C_{\mathrm W}$ and substitute into the second relation of \eqref{eq:weerts-enlarged-relations}.  The $r_2$ and $r_4$ coefficients of $w_2$ simplify using
\[
  C_{\mathrm W}+G_{24}G_{41}G_{12}=D_{\mathrm W},
  \qquad
  G_{24}C_{\mathrm W}+G_{24}G_{41}(G_{14}+G_{12}G_{24})=G_{24}.
\]
For the stated order of $s$, this gives
\begin{equation}
  T_{xs}=\frac{1}{C_{\mathrm W}}
  \begin{bmatrix}
    G_{13}&1&G_{12}&G_{14}+G_{12}G_{24}\\
    G_{24}G_{41}G_{13}&G_{24}G_{41}&1-G_{14}G_{41}&G_{24}
  \end{bmatrix}.
  \label{eq:weerts-measured-map}
\end{equation}
Taking ratios of the entries in \eqref{eq:weerts-measured-map} gives
\begin{align}
  \frac{T_{w_1r_2}}{T_{w_1r_1}}
    &=\frac{G_{12}/C_{\mathrm W}}{1/C_{\mathrm W}}=G_{12},
  &\frac{T_{w_1w_3}}{T_{w_1r_1}}
    &=\frac{G_{13}/C_{\mathrm W}}{1/C_{\mathrm W}}=G_{13},\\
  \frac{T_{w_2r_4}}{T_{w_1r_1}}
    &=\frac{G_{24}/C_{\mathrm W}}{1/C_{\mathrm W}}=G_{24},
  &\frac{T_{w_2r_1}}{T_{w_2r_4}}
    &=\frac{G_{24}G_{41}/C_{\mathrm W}}{G_{24}/C_{\mathrm W}}=G_{41},\\
  \frac{T_{w_1r_4}}{T_{w_1r_1}}-G_{12}G_{24}
    &=G_{14}+G_{12}G_{24}-G_{12}G_{24}=G_{14}.
  \label{eq:weerts-recovery-row1}
\end{align}
Away from the exceptional set where the required denominators vanish, the plant-module fiber is a singleton.  This calculation assumes the stated external channels are available; it does not establish noise-model identifiability or predictor informativity.

Thus both methods measure $w_1,w_2,w_3$.  The reference method identifies the target and composite nuisance dynamics.  The enlarged-subnetwork method identifies the original modules jointly.  It does not reduce node sensors in this example, but it yields a richer identifiable open-loop model.

\subsection{Example 12: parallel path}

Select $G_{ji}$ as the target module.  The node $w_u$ is the intermediate signal on the parallel path.  The schematic network contains
\begin{align}
  w_u&=G_{ui}w_i+\cdots,\\
  w_j&=G_{ji}w_i+G_{ju}w_u+\cdots.
  \label{eq:parallel-path}
\end{align}
Let $w_l=G_{lu}w_u$ be an indirect observation.  Substitution of $w_u=G_{lu}^{-1}w_l$ into \eqref{eq:parallel-path} gives
\begin{equation}
  w_j=G_{ji}w_i+G_{ju}G_{lu}^{-1}w_l+\cdots.
  \label{eq:parallel-abstraction}
\end{equation}
The displayed abstraction preserves $G_{ji}$ as the coefficient of $w_i$.  Its algebraic form alone does not establish a causal predictor or its consistency; the referenced method's assumptions must also hold.

For the simplified displayed plant paths, with omitted terms supplying no additional $w_i$-dependent path and no additional measured transfer channels, including the $u$-equation and measuring only $x=w_j$ gives
\begin{equation}
  T_{w_jw_i}=G_{ji}+G_{ju}G_{ui}.
  \label{eq:parallel-fiber}
\end{equation}
In an unrestricted strictly proper transfer-function class, let $K_{\mathrm p}$ be a nonzero strictly proper perturbation for which the modified modules remain admissible.  In \eqref{eq:parallel-fiber}, replacing $G_{ui}$ by $G_{ui}+K_{\mathrm p}$ and $G_{ji}$ by $G_{ji}-G_{ju}K_{\mathrm p}$ gives
\[
  (G_{ji}-G_{ju}K_{\mathrm p})+G_{ju}(G_{ui}+K_{\mathrm p})
  =G_{ji}+G_{ju}G_{ui}.
\]
When $G_{ju}K_{\mathrm p}\ne0$, the target changes without changing the transfer.  This does not analyze every referenced experiment or a fixed-order parameterization.

Including both displayed equations puts the parallel path inside $\Aset$.  Condition (O3) must still be checked for any feedback paths hidden in the ellipses.  Failure of a departure certificate does not, by itself, prove that enlargement is necessary.

In the simplified topology without omitted feedback or extra paths, if $w_u$ is also measured and an independent excitation $r_u$ enters its equation, then
\[
  T_{w_uw_i}=G_{ui},
  \qquad
  T_{w_jr_u}=G_{ju},
\]
and
\begin{equation}
  G_{ji}
  =T_{w_jw_i}-T_{w_jr_u}T_{w_uw_i}.
  \label{eq:parallel-recovery}
\end{equation}
These additional transfer channels determine the three displayed modules.  They do not establish a minimum sensor requirement or a general measurement advantage for either method.

\subsection{Example 13: loop through the output}

The target is $G_{ji}$, and $w_u$ is the intermediate node on the displayed output loop.  The relevant equations are
\begin{align}
  w_j&=G_{ji}w_i+G_{ju}w_u+\cdots,\\
  w_u&=G_{uj}w_j+\cdots.
  \label{eq:loop-example}
\end{align}
For the simplified displayed loop, with no additional measured channels or omitted $w_i$-dependent paths, eliminating $w_u$ from \eqref{eq:loop-example} gives the measured transfer matrix
\begin{equation}
  T_{w_jw_i}
  =\frac{G_{ji}}{1-G_{ju}G_{uj}}.
  \label{eq:loop-transfer}
\end{equation}
In an unrestricted strictly proper transfer-function class, let $K_{\ell}$ be a nonzero strictly proper perturbation preserving admissibility.  By \eqref{eq:loop-transfer},
\[
  G_{ji}-T_{w_jw_i}G_{ju}K_{\ell}
  =T_{w_jw_i}\bigl[1-G_{ju}(G_{uj}+K_{\ell})\bigr].
\]
Replacing $G_{uj}$ by $G_{uj}+K_{\ell}$ and $G_{ji}$ by the left-hand side preserves the transfer in \eqref{eq:loop-transfer}, provided the modified denominator is nonzero.  The target changes when $T_{w_jw_i}G_{ju}K_{\ell}\ne0$.  This is not an analysis of a fixed-order parameterization or the complete referenced experiment.

The displayed loop is internal when both equations are included in $\Aset$.  Any feedback through the omitted network still requires a separate check of (O3).

An indirect-observation formulation can preserve the target as an individual coefficient under the referenced path-blocking and predictor assumptions.  This schematic alone does not establish that it is preferable, or that measuring $w_u$ suffices for identifiability under every excitation and noise configuration.

\subsection{Example 17: recursive indirect observations}

This example combines a target-parallel path through $w_u$, an indirect observation $w_l$ of $w_u$, a path $w_i\to w_2\to w_l$, and a possible indirect observation $w_3$ of $w_2$.  The referenced method chooses retained signals recursively so that all target-parallel paths and relevant loops are blocked while $G_{ji}$ remains invariant.

For retained predictor-input indices $\Dset_j$, apply the direct informativity condition \eqref{eq:dankers-direct-info} or the two-stage condition \eqref{eq:dankers-twostage-info}, under the referenced predictor assumptions.  Target invariance requires $\breve G_{ji}=G_{ji}$.

No explicit measured transfer map or fiber calculation is supplied here for the corresponding enlarged-subnetwork construction.  Its identifiability therefore remains unassessed; the schematic alone does not justify a positive-dimensional fiber or a minimum measurement requirement.

The recursive schematic does not specify every path through the omitted network, so it does not permit an unconditional orthogonality conclusion for the enlarged-subnetwork construction.  A complete topology, a specified predictor information set, and the disturbance assumptions are needed before (O3) can be checked.

\section{Two constructed examples}

The following constructions exhibit two algebraic measurement configurations.  The first retains exactly the target input and output as internal measurements.  The second recovers the target without measuring either endpoint.  Both include the complete relevant dynamic network, so neither establishes conservatism or minimality of a boundary-enlargement rule.

\subsection{Example A: model enlargement with only the target endpoints measured}

Let $G_a$ be the target module.  Consider the network
\begin{align}
  w_2 &= G_a w_1+G_b w_3+r_2+v_2,
  \label{eq:endpoint-w2}\\
  w_3 &= G_c w_2+r_3+v_3,
  \label{eq:endpoint-w3}
\end{align}
The three modules are unknown, stable, and strictly proper, and the closed-loop network is assumed stable.  Only the target endpoints $w_1,w_2$ are measured node signals; $r_2,r_3$ are measured external excitations.  Here $w_1$ is an exogenous boundary input receiving no feedback from the network.  Assume $(w_1,r_2,r_3)$ has positive-definite joint spectrum and is independent of the disturbances.  For orthogonality, impose diagonal disturbance filters driven by mutually independent innovations; these assumptions are additional to topology and strict properness.

Define the enlarged module set, measurable input vector, ordered module-input vector, and feedback denominator by
\[
  \Aset=\{G_a,G_b,G_c\},
  \qquad
  s=\begin{bmatrix}w_1&r_2&r_3\end{bmatrix}^{\!T},
\]
\[
  u^{\Aset}=\begin{bmatrix}w_1&w_3&w_2\end{bmatrix}^{\!T},
  \qquad D_{\mathrm A}=1-G_bG_c.
\]
The direct MISO configuration for \eqref{eq:endpoint-w2} measures $w_1,w_2,w_3$.  Substituting \eqref{eq:endpoint-w3} into \eqref{eq:endpoint-w2} gives
\begin{equation}
  D_{\mathrm A}w_2=G_aw_1+r_2+G_br_3+v_2+G_bv_3,
  \label{eq:endpoint-immersed}
\end{equation}
so the coefficient multiplying $w_1$ becomes $G_a/(1-G_bG_c)$ rather than $G_a$.

The selected $\Aset$-part is the complete relevant dynamic part of the construction, so no unknown boundary leaves it and returns; condition (O3) is vacuous.  Conditions (O1)--(O2) hold by the strict-properness and disturbance-separation assumptions.

Suppressing disturbances in \eqref{eq:endpoint-immersed} and dividing by $D_{\mathrm A}$ gives
\begin{equation}
  T_{w_2s}
  =\frac{1}{D_{\mathrm A}}\begin{bmatrix}G_a&1&G_b\end{bmatrix}.
  \label{eq:endpoint-map}
\end{equation}
Write the entries in \eqref{eq:endpoint-map} as $(T^{\mathrm A}_1,T^{\mathrm A}_2,T^{\mathrm A}_3)$.  Their ratios give
\[
  \frac{T^{\mathrm A}_1}{T^{\mathrm A}_2}=G_a,
  \qquad
  \frac{T^{\mathrm A}_3}{T^{\mathrm A}_2}=G_b,
  \qquad
  \frac{T^{\mathrm A}_2-1}{T^{\mathrm A}_3}
  =\frac{1-D_{\mathrm A}}{G_b}=G_c.
\]
Thus, for nonzero $G_b$, the recovery equations are
\begin{equation}
  D_{\mathrm A}=\frac{1}{T^{\mathrm A}_2},
  \qquad
  G_a=\frac{T^{\mathrm A}_1}{T^{\mathrm A}_2},
  \qquad
  G_b=\frac{T^{\mathrm A}_3}{T^{\mathrm A}_2},
  \qquad
  G_c=\frac{T^{\mathrm A}_2-1}{T^{\mathrm A}_3}.
  \label{eq:endpoint-recovery}
\end{equation}
Every module is fixed by \eqref{eq:endpoint-recovery}, so the generic fiber is a singleton.  Using \eqref{eq:endpoint-map} and $w_3=G_cw_2+r_3$ from \eqref{eq:endpoint-w3}, the plant-input transfer matrix is
\begin{equation}
  T_{u^{\Aset}s}=\begin{bmatrix}
    1&0&0\\
    G_cG_a/D_{\mathrm A}&G_c/D_{\mathrm A}&1/D_{\mathrm A}\\
    G_a/D_{\mathrm A}&1/D_{\mathrm A}&G_b/D_{\mathrm A}
  \end{bmatrix}.
  \label{eq:endpoint-input-matrix}
\end{equation}
Expanding the determinant of \eqref{eq:endpoint-input-matrix} along its first row gives
\begin{equation}
  \det T_{u^{\Aset}s}
  =\frac{G_cG_b-1}{D_{\mathrm A}^2}=-\frac{1}{D_{\mathrm A}}.
  \label{eq:endpoint-rank}
\end{equation}
The nonzero determinant in \eqref{eq:endpoint-rank} and the assumed positive-definite excitation spectrum satisfy \eqref{eq:informativity}.  Recovery in \eqref{eq:endpoint-recovery} uses one fewer node sensor than the direct MISO configuration.  This does not establish predictor informativity, consistency, or minimum sensor requirements among all abstraction configurations.

\subsection{Example B: neither target endpoint is measured}

Let $G_b$ be the target module.  Consider the network
\begin{align}
  w_1 &= G_a(w_3+r_0)+r_1+v_1,
  \label{eq:constructed-w1}\\
  w_2 &= G_b w_1+G_d w_3+r_2+v_2,
  \label{eq:constructed-w2}\\
  w_3 &= G_c w_2+r_3+v_3,
  \label{eq:constructed-w3}
\end{align}
The four modules are unknown, stable, and strictly proper, and the closed-loop network is assumed stable.  The external excitations are independent of the disturbances and have positive-definite joint spectrum.  For orthogonality, impose diagonal disturbance filters driven by mutually independent innovations.  The enlarged configuration takes
\[
  \Aset=\{G_a,G_b,G_c,G_d\}
\]
but measures only $w_3$, together with
\[
  s=\begin{bmatrix}r_0&r_1&r_2&r_3\end{bmatrix}^{\!T}.
\]
Order the module inputs as
\[
  u^{\Aset}
  =\begin{bmatrix}w_3+r_0&w_1&w_3&w_2\end{bmatrix}^{\!T},
\]
corresponding respectively to $G_a,G_b,G_d,G_c$, and define
\begin{equation}
  \Delta=1-G_c(G_bG_a+G_d).
  \label{eq:constructed-delta}
\end{equation}
The selected $\Aset$-part is again the complete relevant network, so (O3) is vacuous and (O1)--(O2) hold by construction.

Suppress disturbances in \eqref{eq:constructed-w1}--\eqref{eq:constructed-w3}.  Substitute the first relation into the second, then the second into the third:
\begin{align}
  w_2&=(G_bG_a+G_d)w_3+G_bG_ar_0+G_br_1+r_2,\\
  \Delta w_3&=G_cG_bG_ar_0+G_cG_br_1+G_cr_2+r_3.
  \label{eq:constructed-elimination}
\end{align}
Define $T^{\mathrm B}_i=T_{w_3r_i}$ for this configuration.  Dividing \eqref{eq:constructed-elimination} by $\Delta$ gives
\begin{equation}
  \begin{bmatrix}T^{\mathrm B}_0&T^{\mathrm B}_1&T^{\mathrm B}_2&T^{\mathrm B}_3\end{bmatrix}
  =\frac{1}{\Delta}
   \begin{bmatrix}G_cG_bG_a&G_cG_b&G_c&1\end{bmatrix}.
  \label{eq:constructed-map}
\end{equation}
The first three recovery equations come from ratios of the entries in \eqref{eq:constructed-map}.  For the fourth, use \eqref{eq:constructed-delta}:
\[
  \frac{T^{\mathrm B}_3-1}{T^{\mathrm B}_2}
  =\frac{1-\Delta}{G_c}=G_bG_a+G_d.
\]
For nonzero $G_b$ and $G_c$, the recovery equations are
\begin{align}
  G_a&=\frac{T^{\mathrm B}_0}{T^{\mathrm B}_1}
      =\frac{G_cG_bG_a}{G_cG_b},
  &G_b&=\frac{T^{\mathrm B}_1}{T^{\mathrm B}_2}
      =\frac{G_cG_b}{G_c},
  \label{eq:constructed-recovery1}\\
  G_c&=\frac{T^{\mathrm B}_2}{T^{\mathrm B}_3}
      =\frac{G_c/\Delta}{1/\Delta},
  &G_d&=\frac{T^{\mathrm B}_3-1}{T^{\mathrm B}_2}-G_aG_b.
  \label{eq:constructed-recovery2}
\end{align}
Equations \eqref{eq:constructed-recovery1}--\eqref{eq:constructed-recovery2} fix every module, so the generic fiber is a singleton although neither target endpoint is measured.

For the stated ordering of $u^{\Aset}$, the disturbance-free relations \eqref{eq:constructed-w1}--\eqref{eq:constructed-w3} give
\begin{equation}
  s=Q_{\mathrm B}u^{\Aset},
  \qquad
  Q_{\mathrm B}=\begin{bmatrix}
    1&0&-1&0\\
    -G_a&1&0&0\\
    0&-G_b&-G_d&1\\
    0&0&1&-G_c
  \end{bmatrix}.
  \label{eq:constructed-input-matrix}
\end{equation}
Add $G_a$ times the first row to the second row of $Q_{\mathrm B}$, then $G_b$ times the new second row to the third.  These operations preserve its determinant.  Expanding along the first two columns gives
\[
  \det Q_{\mathrm B}
  =\det\begin{bmatrix}-(G_d+G_bG_a)&1\\1&-G_c\end{bmatrix}
  =G_c(G_d+G_bG_a)-1=-\Delta.
\]
By \eqref{eq:constructed-input-matrix}, $T_{u^{\Aset}s}=Q_{\mathrm B}^{-1}$, so
\begin{equation}
  \det T_{u^{\Aset}s}=\frac{1}{\det Q_{\mathrm B}}=-\frac{1}{\Delta}.
  \label{eq:constructed-rank}
\end{equation}
The nonzero determinant in \eqref{eq:constructed-rank} and the assumed excitation spectrum satisfy \eqref{eq:informativity}.  This does not certify predictor informativity.

The displayed direct MISO configuration uses $w_1,w_2,w_3$.  The measured transfer matrix in \eqref{eq:constructed-map} admits algebraic recovery using only the internal sensor at $w_3$ and four external channels.  This comparison does not establish a minimum measurement requirement for the referenced predictor-input selection or abstraction methods, nor does the rational inverse alone establish consistency of a stochastic predictor.

\section{Orthogonality assessment by example}

The orthogonality question for the method of da Mata, Hansson, and Andersen has the following example-by-example answer.  The table distinguishes conditions (O1)--(O2), which are imposed model assumptions, from condition (O3), which is checked from the selected subnetwork and measurements.

\footnotesize
\begin{longtable}{@{}p{0.16\textwidth}p{0.13\textwidth}p{0.14\textwidth}p{0.22\textwidth}p{0.21\textwidth}@{}}
\caption{Conditions from the orthogonality argument checked for the enlarged-subnetwork construction.}
\label{tab:orthogonality-status}\\
\toprule
Example & (O1) Strict properness & (O2) Disturbance separation & (O3) Boundary information & Orthogonality conclusion \\
\midrule
\endfirsthead
\toprule
Example & (O1) Strict properness & (O2) Disturbance separation & (O3) Boundary information & Orthogonality conclusion \\
\midrule
\endhead
\bottomrule
\endfoot
Van den Hof Ex.~1, SISO
& Required assumption; not determined by the displayed topology
& Required assumption; not determined by the displayed topology
& Verified: the outgoing signal is measured or reconstructible as $w_4-w_3$
& Holds for the proposed method when (O1)--(O2) are imposed \\
\addlinespace
Van den Hof Ex.~1, MISO
& Required assumption
& Required assumption
& Verified: the outgoing signal $w_2$ is measured
& Holds for the proposed method when (O1)--(O2) are imposed \\
\addlinespace
Van den Hof Ex.~2
& Required assumption
& Required assumption
& Verified as in the corresponding SISO configurations
& Holds for the proposed method when (O1)--(O2) are imposed \\
\addlinespace
Van den Hof Ex.~3
& Required assumption
& Required assumption
& Verified: $w_2$ is measured
& Holds under (O1)--(O2); the excitation-rank failure is separate \\
\addlinespace
Van den Hof Ex.~4
& Required assumption
& Required assumption
& Verified: $w_2$ is measured
& Holds under (O1)--(O2), without reconstructing $v_3$ \\
\addlinespace
Dankers running example
& Required assumption
& Required assumption
& Not established from the topology specified here
& Not established; complete paths and disturbance locations are required \\
\addlinespace
Weerts four-node example
& Required assumption
& Required assumption
& Verified for the selected modules in \eqref{eq:weerts-four-node}: $w_3$ receives no signal from $\Aset$
& Holds for this specified partition under (O1)--(O2); minimality is not claimed \\
\addlinespace
Weerts Ex.~12
& Required assumption
& Required assumption
& No displayed return remains; omitted paths are not checked
& Holds for the displayed topology under (O1)--(O2); otherwise conditional \\
\addlinespace
Weerts Ex.~13
& Required assumption
& Required assumption
& No displayed return remains; omitted paths are not checked
& Holds for the displayed topology under (O1)--(O2); otherwise conditional \\
\addlinespace
Weerts Ex.~17
& Required assumption
& Required assumption
& Not established from the recursive schematic
& Not established; complete paths and disturbance locations are required \\
\addlinespace
Constructed Example A
& Verified by construction: $G_a,G_b,G_c$ are strictly proper
& Verified by construction: the disturbances satisfy the separation assumptions
& Verified: the enlarged part is the complete relevant dynamic network, so (O3) is vacuous
& All three conditions hold with only the target endpoints measured \\
\addlinespace
Constructed Example B
& Verified by construction: all four plant modules are strictly proper
& Verified by construction: the disturbances satisfy the separation assumptions
& Verified: the enlarged part is the complete relevant network, so (O3) is vacuous
& All three conditions hold with neither target endpoint measured \\
\end{longtable}
\normalsize

Examples A and B include the complete relevant dynamic network and therefore make (O3) vacuous; the disturbance-separation assumptions must also be imposed.  For the displayed Van den Hof configurations and the selected four-node Weerts partition, the preceding arguments check boundary information, not minimality of the selected model or measurements.  The Dankers running example and Weerts Example~17 are not specified sufficiently here to check (O3).  For Weerts Examples~12 and 13, omitted paths remain conditional.  These qualifications are separate from the algebraic recovery calculations.

\section{Consolidated comparison}

\small
\begin{longtable}{@{}p{0.16\textwidth}p{0.17\textwidth}p{0.17\textwidth}p{0.20\textwidth}p{0.16\textwidth}@{}}
\caption{Summary of the example-by-example comparison.}
\label{tab:summary}\\
\toprule
Example & Comparator measurements and model & Enlarged-subnetwork measurements and modules & Boundary information and excitation & Algebraic identifiability conclusion \\
\midrule
\endfirsthead
\toprule
Example & Comparator measurements and model & Enlarged-subnetwork measurements and modules & Boundary information and excitation & Algebraic identifiability conclusion \\
\midrule
\endhead
\bottomrule
\endfoot
Constructed Example A
& Direct MISO: $w_1,w_2,w_3$
& Only target endpoints $w_1,w_2$; enlarged set $G_a,G_b,G_c$
& O1--O2 imposed; O3 vacuous; plant-input determinant $-1/D_{\mathrm A}$
& Singleton fiber; one fewer internal sensor \\
\addlinespace
Constructed Example B
& Direct MISO: $w_1,w_2,w_3$
& Only $w_3$; enlarged set $G_a,G_b,G_c,G_d$
& O1--O2 imposed; O3 vacuous; plant-input determinant $-1/\Delta$
& Singleton fiber; neither target endpoint is measured \\
\addlinespace
Van den Hof Ex.~1, SISO
& Target input and output; one target module
& Same node signals and same target module
& Measured or reconstructible departure signal; scalar excitation test
& Singleton fiber; equivalent \\
\addlinespace
Van den Hof Ex.~1, MISO
& $w_1,w_2,w_3$; $G_{21},G_{23}$
& Same signals and modules
& O3 checked using measured $w_2$; $T_{us}=I_2$
& Singleton fiber \\
\addlinespace
Van den Hof Ex.~2
& Input, output, and $r_1$; projected scalar input
& Same node sensors; direct measured boundary input
& O3 checked; O1--O2 required; scalar excitation conditions
& Singleton fiber \\
\addlinespace
Van den Hof Ex.~3
& $w_1,w_2,w_3,r_1$; two reconstructed inputs
& $w_1,w_2,w_3$; identity plant-input transfer matrix
& O3 verified; O1--O2 required; only the direct plant-input transfer matrix passes the rank test
& Singleton fiber \\
\addlinespace
Van den Hof Ex.~4
& Reconstructed $v_3$ and known $G_{32}$
& Direct boundary inputs $w_1,w_3$
& O3 verified; O1--O2 required; rank two without disturbance reconstruction
& Singleton fiber; less prior module knowledge \\
\addlinespace
Dankers running example
& Retained predictor nodes; invariant target plus immersed nuisance modules
& Boundary nodes plus selected outputs; original modules in enlarged $\Aset$
& O3 is not established for the incompletely specified topology
& Reference target is direct; proposed method requires a fiber calculation \\
\addlinespace
Weerts four-node example
& $w_1,w_2,w_3$; target plus inverse-module nuisance coefficient
& $w_1,w_2,w_3$; five original modules
& O3 checked for the selected partition; no minimum enlargement is established
& Original modules algebraically recoverable with the stated external channels \\
\addlinespace
Weerts Ex.~12
& $w_j,w_i,w_l$; target preserved
& Simplified $w_j,w_i$ configuration absorbs $w_u$
& No displayed return remains after enlargement; omitted paths remain conditional
& Single displayed channel does not fix the target in an unrestricted module class \\
\addlinespace
Weerts Ex.~13
& $w_j,w_i,w_l$; target preserved
& Simplified $w_j,w_i$ configuration absorbs loop node
& No displayed return remains after enlargement; omitted paths remain conditional
& Single displayed channel does not fix the target in an unrestricted module class \\
\addlinespace
Weerts Ex.~17
& Recursive path-blocking retained sets
& Enlarged original-module model with selected outputs
& O3 not established; complete topology and predictor information are required
& Requires a specified measured transfer matrix and identifiability analysis \\
\end{longtable}
\normalsize

\section{Conclusions}

The comparisons support the following conclusions.

\begin{enumerate}
  \item \textbf{The three assessment criteria are separate.}  Conditions (O1)--(O3) support orthogonality in \eqref{eq:orthogonality}.  The plant-input excitation rank test in \eqref{eq:informativity} does not establish predictor informativity.  Plant-module identifiability concerns the fiber in \eqref{eq:fiber}.

  \item \textbf{Nonlinear dependence does not preclude recovery.}  The measured transfer matrix in \eqref{eq:weerts-measured-map} contains products and feedback denominators, but \eqref{eq:weerts-recovery-row1} recovers each original module.

  \item \textbf{Measured boundary signals do not certify orthogonality or informativity.}  Identity or selection columns in $T_{us}$ simplify the plant-input excitation rank test.  They do not establish disturbance independence, (O3), or predictor informativity.

  \item \textbf{The constructed examples reduce measurements algebraically.}  Equation \eqref{eq:endpoint-recovery} recovers the modules with only the target endpoints measured; \eqref{eq:constructed-recovery1}--\eqref{eq:constructed-recovery2} recover them with neither endpoint measured.  These comparisons concern the stated direct MISO configurations, not minimum-sensor implementations of every referenced method.  Consistency still requires the predictor assumptions.

  \item \textbf{Measurement reduction is not universal.}  The Van den Hof configurations retain the same node sensors; the enlarged-subnetwork configurations avoid input projection or disturbance reconstruction in the stated comparisons.

  \item \textbf{Incomplete schematics do not settle a comparison.}  The single-channel transfers in \eqref{eq:parallel-fiber} and \eqref{eq:loop-transfer} do not determine the target in an unrestricted transfer-function class.  They do not establish minimum measurement requirements or results for a complete referenced experiment.

  \item \textbf{Target recovery does not require recovery of every nuisance module.}  The target can be constant on the fiber in \eqref{eq:fiber} even when nuisance modules vary.

  \item \textbf{A generic certificate need not be universal.}  A proved generic certificate of (O3) supports generic orthogonality under the remaining assumptions.  Its admissible parameter class and exceptional set must be specified.  The rank equality in \eqref{eq:boundary-rank-screen} alone does not establish causal reconstruction.  No weaker topology-only algorithm or minimum enlargement is proved here.
\end{enumerate}

\subsection{Research directions for configuration selection}

The comparisons do not yet give a procedure for selecting an enlarged-subnetwork configuration for a single target module.  Such a procedure should select the added original modules, measured node signals, and measured external excitations simultaneously.  Its three costs are the number of added original modules, the number of measured node signals, and the number of independently excited external signals.  A selected configuration is preferable to another when none of these three costs is larger and at least one is smaller.  Reporting all configurations for which no preferable configuration exists avoids imposing an unsupported exchange rate between sensors, excitations, and added modules.

Topology and strict properness alone do not specify the admissible parameter class needed for a generic statement, and they do not imply disturbance separation or excitation.  A topology-based result should therefore first specify a finite-dimensional class of stable strictly proper modules, including order or delay restrictions, together with the disturbance and experiment assumptions.  It should then establish generic conditions separately for causal boundary reconstruction, predictor informativity, and constancy of the target module on the fiber in \eqref{eq:fiber}.  A negative result showing that one of these conditions cannot be decided from topology and strict properness alone would also delimit the attainable guarantee.

The following research tasks would turn the comparison into a selection method:
\begin{enumerate}
  \item Derive graph conditions that reject a candidate configuration when a returning contribution cannot be causally reconstructed from its permitted predictor information.  The conditions must account for delays, because the rational rank screen in \eqref{eq:boundary-rank-screen} does not exclude a noncausal inverse.
  \item Derive topology-based lower bounds on the number of independent excitation directions and measured transfer channels needed to identify the selected original modules.  These bounds can remove candidates before a transfer calculation.
  \item For the remaining candidates, construct the measured transfer matrix and test whether the target module is constant on the fiber in \eqref{eq:fiber}.  The test need not require recovery of every nuisance module.
  \item Search over the remaining candidates and return those for which no configuration has fewer added modules, fewer measured node signals, and fewer independent external excitations simultaneously.
  \item Compare the returned configurations with the predictor-input selections of Van den Hof et al.\ \cite{VandenHof2013} and Dankers et al.\ \cite{Dankers2016}, counting measured node signals, independent external excitations, and additionally identified original modules.
\end{enumerate}

This programme would retain the present separation between orthogonality, informativity, and identifiability.  The constructed examples provide test cases, but they do not establish that their sensor or excitation counts are minimal.

\begin{thebibliography}{99}

\bibitem{VandenHof2013}
P.~M.~J. Van den Hof, A.~G. Dankers, P.~S.~C. Heuberger, and X.~Bombois,
``Identification of dynamic models in complex networks with prediction error methods---basic methods for consistent module estimates,''
\emph{Automatica}, vol.~49, no.~10, pp.~2994--3006, 2013.

\bibitem{Dankers2016}
A.~G. Dankers, P.~M.~J. Van den Hof, X.~Bombois, and P.~S.~C. Heuberger,
``Identification of dynamic models in complex networks with prediction error methods: Predictor input selection,''
\emph{IEEE Transactions on Automatic Control}, vol.~61, no.~4, pp.~937--952, 2016.

\bibitem{Weerts2020}
H.~H.~M. Weerts, J.~Linder, M.~Enqvist, and P.~M.~J. Van den Hof,
``Abstractions of linear dynamic networks for input selection in local module identification,''
\emph{Automatica}, vol.~117, Art.~no.~108975, 2020.

\bibitem{HanssonDaMataAndersen2026}
A.~Hansson, J.~V.~G. da Mata, and M.~S. Andersen,
``Informativity and identifiability for identification of networks of dynamical systems,''
\emph{arXiv preprint arXiv:2602.20107}, 2026.

\bibitem{ZhangFernandoDarouachSFO}
Y.~Zhang, T.~Fernando, and M.~Darouach,
``Functional observability, structural functional observability and optimal sensor placement,''
\emph{arXiv preprint arXiv:2307.08923}, version~3, 2024.

\bibitem{DaMataHanssonAndersenSubnetwork}
J.~V.~G. da Mata, A.~Hansson, and M.~S. Andersen,
``Prediction error method for identification of subnetworks of dynamical systems,''
unpublished manuscript supplied for the present comparison.

\end{thebibliography}

\end{document}
